\documentclass{../shared/classes/utmreport}

\selectlanguage{english}
\utmlogo{../shared/assets/UTM_Logo_English.jpeg}
\institution{Technical University of Moldova}
\faculty{Faculty of Computers, Informatics, and Microelectronics}
\department{Department of Software Engineering}
\studyprogramme{Software Engineering Study Programme}
\documenttype{Internship report}
\reporttitle{Project title}
\team{Team No. 1}
\teammember{Student 1}{First name LAST NAME}{FAF-241}
\teammember{Student 2}{First name LAST NAME}{FAF-241}
\teammember{Student 3}{First name LAST NAME}{FAF-243}
\teammember{Student 4}{First name LAST NAME}{FAF-243}
\internshipmentor{Mentor name from the internship contract}
\submissiondate{day / month / year}
\cityyear{Chișinău}{2026}

\begin{document}
\makeutminternshiptitlepage
\tableofcontents
\clearpage

\begin{utmabstract}[Abstract]
Summarize the problem, the implemented solution, the methods used, and the observed results. Keep completed work separate from planned development. This document is a reusable template; its fields and guidance must be replaced with the team's own report content.
\utmkeywords{domain, requirements, architecture, security, implementation, testing}
\end{utmabstract}

\utmintroduction
Explain the project context, the people who need the system, the problem, the internship objectives, and the scope of the work. Introduce the organization using a source that supports the description. Define abbreviations at first use. Cite adapted ideas, diagrams, tables and technical claims beside the claims they support.

\chapter{Domain analysis}
Explain the domain, existing approaches, user needs and requirements before describing the design.
\section{Problem and user needs}
Describe the people, their workflow and the concrete difficulty being addressed.
\section{Existing approaches and requirements}
Compare relevant approaches using cited evidence and identify the resulting requirements.

\chapter{System design and security}
Explain the design and its decisions before presenting the detailed architecture.
\section{Architecture and technologies}
Cite each reused source at its first relevant use \cite{source}. Introduce each diagram in prose, explain how the components interact, and justify the selected technologies. Put figure captions below figures and table captions above tables.
\section{Security}
Explain risks, access control, data handling, and tested protection mechanisms. Label proposals and untested controls as such.

\chapter{Implementation and testing}
Describe the functionality actually implemented and the evidence for its behavior.
\section{Implementation}
Explain the relevant source code. Code longer than half a page belongs in an appendix referenced here; use Courier New 10 pt with single spacing.
\section{Testing and results}
Describe inputs, expected behavior, observed results and limitations. Illustrative values must not be presented as measurements.

\chapter{Team contributions and future work}
Separate the contribution of each member from the team's shared result.
\section{Contributions and lessons}
State concrete responsibilities, delivered work and lessons. Do not invent attendance, biographies or individual contributions.
\section{Future work}
Describe remaining problems, intended improvements and the next implementation steps.

\utmconclusions
Relate the observed results to the objectives, state the main limitations, and distinguish implemented functionality from planned work.

\begin{thebibliography}{9}
\bibitem{source} Author or organization. \textit{Source title}. Publication date. For an online source, add its URL and actual access date. Each numbered entry must support a citation in the report; arrange references in first-use order.
\end{thebibliography}

\appendix
\chapter{Supplementary source code}
Include longer code listings and supporting material referenced in the main text.

% The accepted Filetracker package used one completed team-level summary last.
% Add separate member summaries only if the course requests them. Leave signatures blank.
\chapter*{Internship Summary Report}
\addcontentsline{toc}{chapter}{Internship Summary Report}
\begin{center}
\textbf{on the activities carried out during the internship}
\end{center}
\noindent Host organisation: \hrulefill\\
Department (unit): \hrulefill\\
Faculty of Computers, Informatics, and Microelectronics\\
Study programme: Software Engineering\\
Mode of study: full time\\
Year of study: \rule{1cm}{0.3pt}\quad Group: \rule{2cm}{0.3pt}\\
Internship period: from \rule{2cm}{0.3pt}\quad to \rule{2cm}{0.3pt}

\noindent\textbf{Activities carried out and results obtained:}

\noindent Activities carried out: \hrulefill\par\vspace{0.3cm}
\noindent Results obtained: \hrulefill\par\vspace{0.3cm}
\noindent\textbf{Difficulties encountered during the internship:}\par
\noindent\hrulefill\par\vspace{0.3cm}
\noindent\textbf{Conclusions and recommendations:}\par
\noindent\hrulefill\par\vspace{0.3cm}
\noindent\hrulefill\par\vspace{0.3cm}
\noindent Date of completion: \rule{0.6cm}{0.3pt} / \rule{0.6cm}{0.3pt} / 20\rule{0.6cm}{0.3pt}\\
Signature of the student: \rule{4cm}{0.3pt}
% Replace blank activity/result/conclusion lines with the student's actual account.
% Assess relevance, practical usefulness, lessons and the host's cooperation.
% Omit difficulties only if they do not apply. Keep the signature blank until signed.
\end{document}
